\section{Problem Formulation and Research Objectives}
\label{sec:problem}

\subsection{Problem Statement}
Let $\mathcal{D}$ be the set of authoritative documents an institution
publishes about itself (web pages, PDF brochures, regulations), and let
$\mathcal{G}=(V,E)$ be a graph model of its physical campus in which
vertices $V$ are navigable points (entrances, junctions, corridors, rooms,
staircases, landmarks) and edges $E$ are walkable connections with
traversal costs. Let $\mathcal{P}$ be a set of 360$^{\circ}$ panoramic
images, each associated with a vertex $v\in V$.

A user issues either a natural-language query $q$ or a spatial request
$(v_s, v_g)$. The system must:

\begin{itemize}
\item \textbf{P1 (grounded answering).} For a query $q$, return an answer
$a$ together with a set of provenance references $S\subseteq\mathcal{D}$
such that every verifiable claim in $a$ is entailed by
$\bigcup_{s\in S}\mathrm{text}(s)$, \emph{or} return an explicit refusal
when $\mathcal{D}$ does not adequately support an answer. The system must
prefer refusal to an unsupported answer.
\item \textbf{P2 (routing).} Assign $q$ to the specialist handler(s) whose
domain it belongs to, deterministically and explainably where possible.
\item \textbf{P3 (wayfinding).} For $(v_s,v_g)$, return a minimum-cost walk
$\pi = (v_s,\dots,v_g)$ in $\mathcal{G}$, its total cost, an estimated
walking time, an accessibility flag, and human-readable turn-by-turn
directions.
\item \textbf{P4 (tour traversal).} Support both manual step-by-step and
automatic sequential traversal of $\mathcal{P}$ along $\mathcal{G}$,
including transitions between floors, while preserving camera-heading
continuity.
\item \textbf{P5 (integration).} Serve P1--P4 from one application over one
shared campus model, deployable on commodity hardware with no external
paid API.
\end{itemize}

\subsection{Design Constraints}
The implemented system operates under constraints that shape every
subsequent design decision:

\begin{enumerate}
\item \emph{No external LLM API.} All generation and embedding runs locally
(Ollama-served Llama~3.2; \code{all-MiniLM-L6-v2}). This bounds model
capability and makes each generation call cost several seconds of CPU time.
\item \emph{Authoritative sources only.} The knowledge base is built
exclusively from the institution's own domain; third-party aggregators,
encyclopedias, and search snippets are excluded by policy.
\item \emph{Groundedness over coverage.} A wrong answer is a more serious
failure than a refusal, because the system's purpose is to be a trustworthy
substitute for the official sources.
\item \emph{Auditability.} Routing decisions and low-confidence refusals
must be explainable after the fact from logs.
\item \emph{Placeholder-to-real asset path.} Panorama imagery and surveyed
coordinates are placeholders; the code must accept real assets by
file/identifier substitution without logic changes.
\end{enumerate}

\subsection{Research Objectives}
\begin{description}
\item[O1] Design and implement a hybrid retrieval pipeline that combines
semantic and exact-term matching for an acronym- and number-heavy
institutional corpus (Section~\ref{sec:rag}).
\item[O2] Design a layered generation-time safety mechanism --- confidence
gating plus post-hoc claim verification --- that reduces the risk of
unsupported institutional claims without a second model call
(Section~\ref{sec:rag}).
\item[O3] Design a routing mechanism that is deterministic and auditable
for covered phrasings and gracefully falls back to a learned classifier for
uncovered ones (Section~\ref{sec:agents}).
\item[O4] Design a scene-graph data structure for 360$^{\circ}$ tour
navigation that makes cross-floor traversal transparent to traversal code
(Section~\ref{sec:tour}).
\item[O5] Implement graph-based wayfinding whose heuristic adapts to
whatever spatial information is available and remains correct in all cases
(Section~\ref{sec:nav}).
\item[O6] Define a rigorous evaluation protocol for all of the above, and
report honestly the measurements that currently exist
(Sections~\ref{sec:eval}--\ref{sec:results}).
\end{description}

\subsection{Non-Goals}
The following are explicitly out of scope for the implemented artifact:
real-time indoor localization; outdoor GPS turn-by-turn guidance (a
prototype was built and then removed); automatic translation of retrieved
evidence (the assistant answers in the user's selected language via the
LLM, but retrieval remains English); authentication and multi-tenant
isolation; and streaming token-by-token responses.
